\sectionkicker{Source-backed technical notes}
\subsection*{Serving Stackは同じ月に何を吸収したか: Technical Notes}
本欄は記事本文で圧縮した一次資料上の情報を、比較・再検証しやすい形で整理したものである。時系列、確認済みの主張、留意点、一次資料URLを掲載する。
\medskip
\subsection*{Theme at a glance}
\addcontentsline{toc}{subsection}{Theme at a glance}
\begin{center}
\footnotesize
\begin{tabularx}{\linewidth}{@{}p{0.20\linewidth}p{0.10\linewidth}p{0.14\linewidth}X@{}}
\toprule
資料 & 位置づけ & 種別 & 時系列 \\
\midrule
vLLM v0.15.1 / v0.16.0 & 主要資料 & Framework & 2026-02-04 (プロジェクト公開); 2026-02-25 (プロジェクト公開) \\
SGLang v0.5.9 & 補足資料 & Framework & 2026-02-24 (プロジェクト公開) \\
Transformers v5.1 / v5.2 model-family integrations & 補足資料 & Framework & 2026-02-05 (プロジェクト公開); 2026-02-16 (プロジェクト公開) \\
\bottomrule
\end{tabularx}
\end{center}
\normalsize

\begin{technicalnote}{vLLM v0.15.1 / v0.16.0}{主要資料}
\begingroup
% reader-facing Technical Notes break policy
\widowpenalty=10000
\clubpenalty=10000
\displaywidowpenalty=10000
\begin{tabularx}{\linewidth}{@{}>{\bfseries}p{0.22\linewidth}X@{}}
組織 & vLLM Project \\
種別 & Framework \\
時系列 & 2026-02-04 (プロジェクト公開); 2026-02-25 (プロジェクト公開) \\
\end{tabularx}
\smallskip
{\bfseries 一次資料から整理したtechnical points}
\begin{itemize}[leftmargin=1.5em,itemsep=0.35em]
\item \textbf{一次情報で確認できる事実}: vLLM v0.16.0は2026年2月25日に公開され、release noteではbranch cutが2026年2月8日だったと記載されている。
\item \textbf{一次情報で確認できる事実}: v0.16.0では、streaming-audio/Voxtral realtime対応を持つWebSocket-based Realtime API、Responses API改善、tool-calling修正、unified parallel drafting、大規模なXPU変更が追加された。
\item \textbf{Project claim}: releaseではasync schedulingとpipeline parallelismによりend-to-end throughputが30.8\%、TPOTが31.8\%改善したと報告されているが、これらはproject-reported measurementである。
\item \textbf{一次情報で確認できる事実}: v0.15.1は2026年2月4日に公開され月初のcontextとなるが、2月のsystems signalとしてはv0.16.0の方が強い。
\end{itemize}
\begin{minipage}{\linewidth}
% reader-facing Technical Notes generic boundary/source tail group
{\bfseries 読む際の境界}
\begin{itemize}[leftmargin=1.5em,itemsep=0.35em]
\item \textbf{分析上の留意点}: v0.16.0のbranch-cut日を考慮すると、2月後半のmodel changeがすべて既に統合済みだったことのEvidenceにはできない。
\item \textbf{分析上の留意点}: performance値はproject release noteのmeasurementであり、独立再現していない。
\item \textbf{分析上の留意点}: このseriesはpatch releaseとmajor releaseをまとめているため、編集上はv0.16.0を中心に扱う。
\end{itemize}
\begin{samepage}
{\bfseries 一次資料}
\begingroup\sloppy
\Urlmuskip=0mu plus 2mu\relax
\begin{itemize}[leftmargin=1.5em,itemsep=0.25em]
\item {\scriptsize\url{https://github.com/vllm-project/vllm/releases/tag/v0.15.1}}
\item {\scriptsize\url{https://github.com/vllm-project/vllm/releases/tag/v0.16.0}}
\end{itemize}
\endgroup
\end{samepage}
\end{minipage}
\endgroup
\end{technicalnote}
\medskip

\begin{technicalnote}{SGLang v0.5.9}{補足資料}
\begingroup
% reader-facing Technical Notes break policy
\widowpenalty=10000
\clubpenalty=10000
\displaywidowpenalty=10000
\begin{tabularx}{\linewidth}{@{}>{\bfseries}p{0.22\linewidth}X@{}}
組織 & SGLang Project \\
種別 & Framework \\
時系列 & 2026-02-24 (プロジェクト公開) \\
\end{tabularx}
\smallskip
{\bfseries 一次資料から整理したtechnical points}
\begin{itemize}[leftmargin=1.5em,itemsep=0.35em]
\item \textbf{一次情報で確認できる事実}: SGLang v0.5.9は2026年2月24日に公開され、Kimi-K2.5、GLM-5、Qwen3.5、MiniMax M2.5、および複数のmultimodal/diffusionモデルファミリへの対応が追加された。
\item \textbf{一次情報で確認できる事実}: このreleaseでは、Anthropic互換API endpointに加え、inference、MoE、multimodal、diffusion runtimeに関する大幅な変更が追加された。
\item \textbf{Project claim}: overlapped LoRA loadingでTTFTを約78\%低減、Blackwell上でDeepSeek V3.2 NSAを3〜5倍高速化したとする数値は、独立測定ではなくproject側の報告である。
\end{itemize}
\begin{minipage}{\linewidth}
% reader-facing Technical Notes generic boundary/source tail group
{\bfseries 読む際の境界}
\begin{itemize}[leftmargin=1.5em,itemsep=0.35em]
\item \textbf{分析上の留意点}: 性能値はprojectのrelease noteに基づくもので、本調査では独立に再現していない。
\item \textbf{分析上の留意点}: support追加は実装が利用可能になったことを示すが、model family間の品質同等性やbenchmarkの妥当性を証明するものではない。
\end{itemize}
\begin{samepage}
{\bfseries 一次資料}
\begingroup\sloppy
\Urlmuskip=0mu plus 2mu\relax
\begin{itemize}[leftmargin=1.5em,itemsep=0.25em]
\item {\scriptsize\url{https://github.com/sgl-project/sglang/releases/tag/v0.5.9}}
\end{itemize}
\endgroup
\end{samepage}
\end{minipage}
\endgroup
\end{technicalnote}
\medskip

\begin{technicalnote}{Transformers v5.1 / v5.2 model-family integrations}{補足資料}
\begingroup
% reader-facing Technical Notes break policy
\widowpenalty=10000
\clubpenalty=10000
\displaywidowpenalty=10000
\begin{tabularx}{\linewidth}{@{}>{\bfseries}p{0.22\linewidth}X@{}}
組織 & Hugging Face \\
種別 & Framework \\
時系列 & 2026-02-05 (プロジェクト公開); 2026-02-16 (プロジェクト公開) \\
\end{tabularx}
\smallskip
{\bfseries 一次資料から整理したtechnical points}
\begin{itemize}[leftmargin=1.5em,itemsep=0.35em]
\item \textbf{一次情報で確認できる事実}: Transformers v5.2.0は2026年2月16日に公開され、VoxtralRealtime、GLM-5、Qwen3.5/Qwen3.5 MoE、VibeVoice Acoustic Tokenizerへの明示的なsupportが追加された。
\item \textbf{Project claim}: VoxtralRealtime implementationでは、audio encoderとMistral-based language-model decoderにtime conditioningとcausal-convolution cacheを組み合わせたstreaming ASRが説明されている。
\item \textbf{一次情報で確認できる事実}: Transformers v5.1.0は2月前半に公開され、GLM-OCRなど追加のmodel-family integrationを含んでいた。
\item \textbf{分析上の留意点}: Transformers release note内に引用されたvendorのperformance記述は、Hugging Faceによる独立検証ではない。
\end{itemize}
\begin{minipage}{\linewidth}
% reader-facing Technical Notes generic boundary/source tail group
{\bfseries 読む際の境界}
\begin{itemize}[leftmargin=1.5em,itemsep=0.35em]
\item \textbf{分析上の留意点}: framework integrationはTransformersへimplementation/supportが入ったことを示すが、model vendorのbenchmark claimを独立検証するものではない。
\item \textbf{分析上の留意点}: v5.1.0はsupporting contextであり、2月のintegration milestoneとしてはv5.2.0の方が強い。
\item \textbf{分析上の留意点}: v5.2 releaseにはbreaking interface changeも含まれるため、model supportを単純なcompatibility listだけに還元すべきではない。
\end{itemize}
\begin{samepage}
{\bfseries 一次資料}
\begingroup\sloppy
\Urlmuskip=0mu plus 2mu\relax
\begin{itemize}[leftmargin=1.5em,itemsep=0.25em]
\item {\scriptsize\url{https://github.com/huggingface/transformers/releases/tag/v5.1.0}}
\item {\scriptsize\url{https://github.com/huggingface/transformers/releases/tag/v5.2.0}}
\end{itemize}
\endgroup
\end{samepage}
\end{minipage}
\endgroup
\end{technicalnote}
\medskip
